% Generated draft: requires mathematical and source review.
\documentclass[11pt]{article}
\usepackage[margin=25mm]{geometry}
\usepackage{fontspec}
\setmainfont{Latin Modern Roman}
\usepackage{amsmath,amssymb,mathtools,unicode-math}
\setmathfont{Latin Modern Math}
\usepackage{enumitem,needspace,xurl,hyperref}
\setlength{\parindent}{0pt}
\setlength{\parskip}{5pt}
\setlength{\emergencystretch}{3em}
\newcommand{\N}{\mathbb{N}}
\newcommand{\Z}{\mathbb{Z}}
\newcommand{\Q}{\mathbb{Q}}
\newcommand{\R}{\mathbb{R}}
\newcommand{\C}{\mathbb{C}}
\newcommand{\F}{\mathbb{F}}
\newcommand{\A}{\mathbb{A}}
\newcommand{\PP}{\mathbb{P}}
\newcommand{\E}{\mathbb{E}}
\newcommand{\eps}{\varepsilon}
\newcommand{\dd}{\,\mathrm{d}}
\begin{document}
\section*{UnsolvedMath ID: 30002848}
\textbf{Source problem code:} catalogue entry
\subsection{Definitions and mathematical statement}
For terminology, a nonzero univariate polynomial \(p\in\R[x]\) is \emph{real-rooted} if every complex zero of \(p\) is real; nonzero constant polynomials satisfy this definition vacuously. This terminology does not specify the admissible polynomial class intended by the source, which leaves that class and the treatment of the zero polynomial unspecified.

The supplied record poses two research questions:
\begin{enumerate}
\item Can one define a quantitative measure of departure from real-rootedness and extend inequalities valid under real-rootedness to polynomials outside that class, with a penalty controlled by this measure?
\item Can interlacing-polynomial methods be extended to obtain simultaneous bounds on two eigenvalues?
\end{enumerate}
The supplied background also mentions improving paving block counts as a motivation for the second question, but provides no paving formulation or quantitative target.

These questions are not specified sufficiently to give a unique quantified conjecture. For the first, the source does not identify the polynomial class, the inequalities to be extended, the properties required of the quantitative measure, or the dependence and limiting behavior of the penalty. For the second, it does not identify the matrix or operator class, the two eigenvalues, the hypotheses of the interlacing construction, or the desired simultaneous bounds. These choices remain unresolved; no particular distance, inequality, or spectral estimate is asserted here.
\subsection{Short English statement}
The problem asks for quantitative extensions of inequalities associated with real-rooted polynomials and for interlacing methods that bound two eigenvalues simultaneously, with improved paving block counts mentioned as motivation. The supplied record leaves the hypotheses and quantitative targets unspecified, so this draft preserves the research questions and identifies the missing specifications.
\subsection{Sources}
\begin{itemize}
\item Local UnsolvedMath registry: \texttt{data/unsolvedmath/problems.json}, supplied record with permanent ID \(30002848\), source problem code \texttt{catalogue entry}, title \emph{Quantitative Real-Rootedness and Two-Eigenvalue Interlacing}. No individual proposer is identified; the supplied proposed year is 2015.
\item Supplied report citation: \emph{Arbeitsgemeinschaft: The Kadison-Singer Conjecture (2016)}, Oberwolfach Reports, pp.~927--969, DOI \texttt{10.4171/owr/2015/17}. Supplied source URL and literature source labeled ``original OWR report'': \url{https://doi.org/10.4171/owr/2015/17}. The citation's parenthetical year 2016 and the supplied proposed year 2015 are retained without reconciliation.
\item Supplied literature assessment dated \texttt{2026-08-21}: the record labels the problem open and describes the extract as an unspecific quantitative strengthening concerning two eigenvalues and improved paving block counts, not resolved by the cited Kadison--Singer theorem. This is supplied dated curation, not an independently verified literature assessment.
\item The supplied statement status is \texttt{corrected\_verified}. Its verification note reports that the official Oberwolfach passage was checked directly and that the question, definitions, and notation were restored without strengthening the claim. This draft records that supplied verification claim; it does not claim independent inspection of the report.
\end{itemize}
\end{document}
